\documentclass[a4paper,12pt]{article}

\usepackage{geometry}
\usepackage{cite}
\usepackage{url}
\usepackage{amsfonts}
\usepackage{amsmath}
\geometry{margin=1.5in}
\usepackage{enumitem}

\title{Efficient instruction-level parallelism analysis using tropical semiring \lq{matrix multiplication\rq}.}
\author{Bryan Abi Karam}
\date{\today}
\begin{document}

\maketitle

\section{Introduction}

% WHAT IS ILP?

Instruction-level parallelism (ILP) is a technique commonly
used by high-performance processors to improve performance.
Previous studies \cite{11096386} have investigated the theoretical
limits of ILP in programs using models of their dynamic execution.

% WHAT IS ONE UPPER BOUND FOR PARALLELISM? CRITICAL PATH

An upper bound for parallelism in a program arises from its critical
path: the longest chain of dependent instructions in the program's
dynamic execution. For a machine that takes at least one cycle per
instruction, the program cannot be run faster than the time it takes
to run the critical path sequentially.

Therefore an upper bound for ILP in a program can be calculated as the
number of dynamic instructions in the program's execution divided by the
length of the critical path.

% HOW TO COMPUTE LENGTH OF CRITICAL PATH?

The dependencies of instructions in a program execution can be represented
by its full dynamic dependency graph (DDG) \cite{753330}, which is a directed
acyclic graph with one node per dynamically executed instruction in a program.
The edges of the graph represent dependencies between dynamically executed instructions,
and so the length of the critical path is the length of the longest path through
the program's DDG.

% WHAT IS TROPICAL ANALYSIS?

Tropical analysis is a branch of pure mathematics using the tropical semiring \cite{Chadderz},
also known as the min/max-plus algebra, which is the semiring
\[(\mathbb{R} \cup \{+\infty\}, \oplus, \otimes)\]
with operations
\[
\begin{aligned}
  a \oplus b & = \min(a, b) \text{ or} \max(a, b) \\
  a \otimes b & = a + b
\end{aligned}
\]
The problem of finding the longest path in a DAG can be expressed as a series
of \lq{matrix multiplications\rq} in the max tropical semiring, as shown by
Section I-B of Chadwick \cite{Chadderz}.

% WHAT WILL I DO IN THIS PROJECT?

In this project I will create a library for performing such matrix
calculations efficiently, considering aspects such as vectorisation,
tiling, and representation sparsity. I will work with CPU based compute
and may work with GPU based compute as an extension.

This introduction is adapted from the abstract of Chadwick \cite{Chadderz}.

\section{The project}

\subsection{Project core}
First, a CPU-based library for performing individual matrix computations in the tropical
semiring should be implemented, using a dense representation of matrices.
Techniques such as vectorisation and tiling should be considered for performance gains.

Matrix multiplication expressions will be generated from a program execution's dependency
structure as a graph, with one node per register per dynamic instruction,
in the column-based representation proposed in Section I-B of Chadwick \cite{Chadderz}.

The aforementioned matrix computation library will be used to compute the results
of these expressions, using techniques such as caching and evaluation order optimisation
to maximise efficiency.

Finally, a benchmark suite for testing performance of implementations will be implemented,
and a baseline performance will be established using a \lq{classical\rq} technique for computing
dependency chains (i.e. simulating one instruction at a time naively, computing its dependencies
one-by-one).

\subsection{Project extensions}

\subsubsection{Sparse matrix representation}
The core of the project implements a matrix computation library using a dense
representation of matrices. In an effort to improve memory usage, and potentially
improve performance (as the cost of sparse multiplication is proportional to the
number of non-zeroes in the matrix, rather than just the number of rows and columns),
I will implement a similar CPU-based matrix computation library using a sparse
representation of matrices as an extension to the project core.

The sparse library will be benchmarked against the dense library for register-only
program execution dependency graphs.

\subsubsection{GPU compute}
The core of the project implements a CPU-based matrix computation library.
To exploit the excellent parallel power of GPUs, I will implement two GPU-based matrix
computation libraries, for dense and sparse representations of matrices, using a
programming model such as OpenCL or CUDA.

The GPU libraries will be benchmarked against their CPU counterparts for register-only
program execution dependency graphs.

\subsubsection{Through-memory dependencies}
So far we have only computed results of matrix expressions generated from register-only
program execution dependency graphs. However, for real programs, the true critical path
often involves memory operations.

A simple approach is to naively include memory locations as graph nodes alongside
registers for each dynamic instruction, but this results in a catastrophic explosion of
matrix sizes. Nevertheless, these matrices are extremely sparse, as most memory operations
access only two or three locations at a time, and so are usable with sparse representations.
For this approach, I will benchmark the performance of the two previously implemented sparse
matrix computation libraries (running on CPU and GPU) against the baseline.

An alternative approach is to model the effects of memory operations using tropical semiring
\textit{addition} rather than just multiplication, as described in Section II-B of Chadwick \cite{Chadderz}.
For this approach, I will benchmark the performance of the two previously implemented dense
matrix computation libraries (running on CPU and GPU) against the baseline.

\section{Starting point}
Existing standard libraries for matrix multiplication cannot be used as they are
specialised for linear algebra (as opposed to tropical algebra).
Additionally, techniques such as Strassen's algorithm do not apply in a semiring.

A suite of tools will be provided by Alexandra Chadwick, the project supervisor,
to generate matrix expressions from program executions.
The tools will do the following:
\begin{itemize}
  \item Compute the static dependencies of every static instruction in a program.
  \item Turn the above into a matrix per instruction.
  \item Dump the dynamic flow of execution of a program.
  \item Use the dynamic flow and the per-instruction matrices to create a matrix expression.
\end{itemize}
Generated matrix expressions will be used as input to the matrix computation libraries.

\section{Evaluation plan}
% Outline how you will evaluate the project to determine its success.
% Include criteria, methods, and metrics.
TBC

\section{Plan of work}
TBC

\section{Resource declaration}
% List resources required for the project (hardware, software, data, etc.)
TBC

\bibliography{bibliography}{}
\bibliographystyle{siam}
\end{document}
